% Section owner: Manas (integration)
\section{Introduction}
\label{sec:p2intro}
The project asks whether a CNN separates visually similar clothing categories better than fixed image
descriptors. Phase 1 answered this for one small CNN: three convolution layers followed by two dense layers reached
\PhaseOneCnnAcc{} accuracy on the Fashion-MNIST test set~\cite{xiao2017fashion}, against \PhaseOneSvmAcc{} for the
best classical model, an RBF-SVM on pixels. That result had stated limits. The CNN was trained once, with one
seed, without augmentation or tuning, and its recall on the Shirt class was only \PhaseOneShirt.

Phase 2 keeps the data split and the evaluation protocol of Phase 1 and asks four questions.
\begin{enumerate}
\item \emph{Architecture.} Does a deeper network with batch normalisation, residual blocks and global average
pooling separate the confusable classes better, and is the gain due to the architecture or to the training recipe?
\item \emph{Regularisation and data.} Which of augmentation, dropout, weight decay, label smoothing and the
learning-rate schedule contribute to validation performance?
\item \emph{Hybrid.} Do handcrafted HOG features add information to the learned CNN embedding when both are given
to a classical classifier?
\item \emph{Validation.} Are the differences larger than the variation between training runs and the sampling
noise of the test set, and what do they cost in calibration, robustness and inference time?
\end{enumerate}

The contributions are: a configurable network family in which each architectural claim is tested by switching one
option (Section~\ref{sec:p2arch}); \NRuns{} training runs covering a four-step architecture roadmap and thirteen
one-factor ablations with three seeds each (Sections~\ref{sec:p2abl} and~\ref{sec:p2results}); a feature-level
hybrid with two controls and a confidence-gated cascade (Section~\ref{sec:p2hybrid}); and a validation protocol
with a freeze file, a one-shot test gate, multi-seed statistics, paired significance tests, calibration and a
corruption stress test (Section~\ref{sec:p2protocol}). Results that did not go in the expected direction are
reported as they were measured.
